% OpenStax Calculus Volume 3 -- Section 7.3 Exercises: Applications
% Exercises reproduced from OpenStax, licensed under CC BY 4.0.
% Source: https://openstax.org/books/calculus-volume-3/pages/7-3-applications
\documentclass{exercises}
\title{OpenStax Calculus Volume 3}
\subtitle{Section 7.3 Exercises: Applications}
\sourceurl{https://openstax.org/books/calculus-volume-3/pages/7-3-applications}
\author{}
\date{}

\begin{document}
\maketitle

\begin{enumerate}[start=1]
\item A mass weighing 4 lb stretches a spring 8 in. Find the equation of motion if the spring is released from the equilibrium position with a downward velocity of 12 ft/sec. What is the period and frequency of the motion?
\item A mass weighing 2 lb stretches a spring 2 ft. Find the equation of motion if the spring is released from 2 in. below the equilibrium position with an upward velocity of 8 ft/sec. What is the period and frequency of the motion?
\item A 100-g mass stretches a spring 0.1 m. Find the equation of motion of the mass if it is released from rest from a position 20 cm below the equilibrium position. What is the frequency of this motion?
\item A 400-g mass stretches a spring 5 cm. Find the equation of motion of the mass if it is released from rest from a position 15 cm below the equilibrium position. What is the frequency of this motion?
\item A block has a mass of 9 kg and is attached to a vertical spring with a spring constant of 0.25 N/m. The block is stretched 0.75 m below its equilibrium position and released.
\begin{enumerate}[label=(\alph*)]
  \item Find the position function $x(t)$ of the block.
  \item Find the period and frequency of the vibration.
  \item Sketch a graph of $x(t)$.
  \item At what time does the block first pass through the equilibrium position?
\end{enumerate}
\item A block has a mass of 5 kg and is attached to a vertical spring with a spring constant of 20 N/m. The block is released from the equilibrium position with a downward velocity of 10 m/sec.
\begin{enumerate}[label=(\alph*)]
  \item Find the position function $x(t)$ of the block.
  \item Find the period and frequency of the vibration.
  \item Sketch a graph of $x(t)$.
  \item At what time does the block first pass through the equilibrium position?
\end{enumerate}
\item A 1-kg mass is attached to a vertical spring with a spring constant of 21 N/m. The resistance in the spring-mass system is equal to 10 times the instantaneous velocity of the mass.
\begin{enumerate}[label=(\alph*)]
  \item Find the equation of motion if the mass is released from a position 2 m below its equilibrium position with a downward velocity of 2 m/sec.
  \item Graph the solution and determine whether the motion is overdamped, critically damped, or underdamped.
\end{enumerate}
\item An 800-lb weight (25 slugs) is attached to a vertical spring with a spring constant of 226 lb/ft. The system is immersed in a medium that imparts a damping force equal to 10 times the instantaneous velocity of the mass.
\begin{enumerate}[label=(\alph*)]
  \item Find the equation of motion if it is released from a position 20 ft below its equilibrium position with a downward velocity of 41 ft/sec.
  \item Graph the solution and determine whether the motion is overdamped, critically damped, or underdamped.
\end{enumerate}
\item A 9-kg mass is attached to a vertical spring with a spring constant of 16 N/m. The system is immersed in a medium that imparts a damping force equal to 24 times the instantaneous velocity of the mass.
\begin{enumerate}[label=(\alph*)]
  \item Find the equation of motion if it is released from its equilibrium position with an upward velocity of 4 m/sec.
  \item Graph the solution and determine whether the motion is overdamped, critically damped, or underdamped.
\end{enumerate}
\item A 1-kg mass stretches a spring 61.25 cm. The resistance in the spring-mass system is equal to eight times the instantaneous velocity of the mass.
\begin{enumerate}[label=(\alph*)]
  \item Find the equation of motion if the mass is released from a position 5 m below its equilibrium position with an upward velocity of 10 m/sec.
  \item Determine whether the motion is overdamped, critically damped, or underdamped.
\end{enumerate}
\item A 32-lb weight (1 slug) stretches a vertical spring 128 in. The resistance in the spring-mass system is equal to four times the instantaneous velocity of the mass.
\begin{enumerate}[label=(\alph*)]
  \item Find the equation of motion if it is released from its equilibrium position with a downward velocity of 12 ft/sec.
  \item Determine whether the motion is overdamped, critically damped, or underdamped.
\end{enumerate}
\item A 64-lb weight is attached to a vertical spring with a spring constant of 4.625 lb/ft. The resistance in the spring-mass system is equal to the instantaneous velocity. The weight is set in motion from a position 1 ft below its equilibrium position with an upward velocity of 2 ft/sec. Is the mass above or below the equilibrium position at the end of $\pi$ sec? By what distance?
\item A mass that weighs 8 lb stretches a spring 6 inches. The system is acted on by an external force of $8\sin(8t)\,\text{lb}$. If the mass is pulled down 3 inches and then released, determine the position of the mass at any time.
\item A mass that weighs 6 lb stretches a spring 3 in. The system is acted on by an external force of $8\sin(4t)\,\text{lb}$. If the mass is pulled down 1 inch and then released, determine the position of the mass at any time.
\item Find the charge on the capacitor in an RLC series circuit where $L=40\,\text{H}$, $R=30\,\Omega$, $C=1/200\,\text{F}$, and $E(t)=200\,\text{V}$. Assume the initial charge on the capacitor is $7\,\text{C}$ and the initial current is $0\,\text{A}$.
\item Find the charge on the capacitor in an RLC series circuit where $L=2\,\text{H}$, $R=24\,\Omega$, $C=0.005\,\text{F}$, and $E(t)=12\sin(10t)\,\text{V}$. Assume the initial charge on the capacitor is $0.001\,\text{C}$ and the initial current is $0\,\text{A}$.
\item A series circuit consists of a device where $L=1\,\text{H}$, $R=20\,\Omega$, $C=0.002\,\text{F}$, and $E(t)=12\,\text{V}$. If the initial charge and current are both zero, find the charge and current at time $t$.
\item A series circuit consists of a device where $L=\frac{1}{2}\,\text{H}$, $R=10\,\Omega$, $C=\frac{1}{50}\,\text{F}$, and $E(t)=250\,\text{V}$. If the initial charge on the capacitor is $0\,\text{C}$ and the initial current is $18\,\text{A}$, find the charge and current at time $t$.
\end{enumerate}

\end{document}
